% Options for packages loaded elsewhere
\PassOptionsToPackage{unicode}{hyperref}
\PassOptionsToPackage{hyphens}{url}
\documentclass[
  ignorenonframetext,
]{beamer}
\newif\ifbibliography
\usepackage{pgfpages}
\setbeamertemplate{caption}[numbered]
\setbeamertemplate{caption label separator}{: }
\setbeamercolor{caption name}{fg=normal text.fg}
\beamertemplatenavigationsymbolsempty
% remove section numbering
\setbeamertemplate{part page}{
  \centering
  \begin{beamercolorbox}[sep=16pt,center]{part title}
    \usebeamerfont{part title}\insertpart\par
  \end{beamercolorbox}
}
\setbeamertemplate{section page}{
  \centering
  \begin{beamercolorbox}[sep=12pt,center]{section title}
    \usebeamerfont{section title}\insertsection\par
  \end{beamercolorbox}
}
\setbeamertemplate{subsection page}{
  \centering
  \begin{beamercolorbox}[sep=8pt,center]{subsection title}
    \usebeamerfont{subsection title}\insertsubsection\par
  \end{beamercolorbox}
}
% Prevent slide breaks in the middle of a paragraph
\widowpenalties 1 10000
\raggedbottom
\AtBeginPart{
  \frame{\partpage}
}
\AtBeginSection{
  \ifbibliography
  \else
    \frame{\sectionpage}
  \fi
}
\AtBeginSubsection{
  \frame{\subsectionpage}
}
\usepackage{iftex}
\ifPDFTeX
  \usepackage[T1]{fontenc}
  \usepackage[utf8]{inputenc}
  \usepackage{textcomp} % provide euro and other symbols
\else % if luatex or xetex
  \usepackage{unicode-math} % this also loads fontspec
  \defaultfontfeatures{Scale=MatchLowercase}
  \defaultfontfeatures[\rmfamily]{Ligatures=TeX,Scale=1}
\fi
\usepackage{lmodern}
\ifPDFTeX\else
  % xetex/luatex font selection
\fi
% Use upquote if available, for straight quotes in verbatim environments
\IfFileExists{upquote.sty}{\usepackage{upquote}}{}
\IfFileExists{microtype.sty}{% use microtype if available
  \usepackage[]{microtype}
  \UseMicrotypeSet[protrusion]{basicmath} % disable protrusion for tt fonts
}{}
\makeatletter
\@ifundefined{KOMAClassName}{% if non-KOMA class
  \IfFileExists{parskip.sty}{%
    \usepackage{parskip}
  }{% else
    \setlength{\parindent}{0pt}
    \setlength{\parskip}{6pt plus 2pt minus 1pt}}
}{% if KOMA class
  \KOMAoptions{parskip=half}}
\makeatother
\setlength{\emergencystretch}{3em} % prevent overfull lines
\providecommand{\tightlist}{%
  \setlength{\itemsep}{0pt}\setlength{\parskip}{0pt}}
\usepackage{bookmark}
\IfFileExists{xurl.sty}{\usepackage{xurl}}{} % add URL line breaks if available
\urlstyle{same}
\hypersetup{
  hidelinks,
  pdfcreator={LaTeX via pandoc}}

\author{\texorpdfstring{}{}}
\date{}

\begin{document}

\section{Introduction: The Teleological Attack
Surface}\label{introduction-the-teleological-attack-surface}

\begin{frame}{Series Context}
\protect\phantomsection\label{series-context}
This paper is Part 4 of the \emph{Cognitive Security for Multiagent
Operators} series. Paper 1 \cite{friedman2026cogsec1} established the
Cognitive Integrity Framework (CIF): a formal model of agent cognitive
states
\(\sigma_i = \langle \mathcal{B}_i, \mathcal{G}_i, \mathcal{I}_i, \mathcal{H}_i \rangle\),
a trust calculus with delegation decay, a five-tier adversary taxonomy
(\(\Omega_1\)--\(\Omega_5\)), and five canonical defense mechanisms with
composition algebra. Paper 2 \cite{friedman2026cogsec2} validated these
mechanisms computationally, demonstrating 99.5\% detection across 950
attack scenarios with the recommended defense stack. Paper 3
\cite{friedman2026cogsec3} grounded the framework in biological analogy,
identifying eusocial insect colonies as evolutionary existence proofs
for CIF-like defense architectures.

This paper addresses the remaining question: \textbf{does CIF work in
practice?} We apply the framework across ten critical domains---from
millisecond drone swarm decisions to year-scale diplomatic
deliberations---demonstrating both its universality and the novel
defense patterns that emerge from domain-specific application.
\end{frame}

\begin{frame}{The Ontological Crisis in AI}
\protect\phantomsection\label{the-ontological-crisis-in-ai}
The vulnerability of modern Artificial Intelligence has shifted from the
\emph{epistemic} (what the agent knows) to the \emph{teleological} (what
the agent wants) \cite{waltzman2017weaponization, aiagentssurvey2025}.
\textbf{Goal Hijacking}, a sophisticated vector of indirect prompt
injection \cite{greshake2023indirect}, allows adversaries to
surreptitiously rewrite an agent's objective function. This represents
an ontological crisis for autonomous systems: if an agent cannot trust
the integrity of its own goals, it cannot trust any action it
calculates.

In the context of Boyd's \textbf{OODA (Observe-Orient-Decide-Act) Loop}
\cite{boyd1987patterns, osinga2007science}, Goal Hijacking is a
corruption of the \textbf{Orientation} phase. The agent correctly
Observes the world, but its internal Orientation---the synthesis of
heritage, culture, and genetic code (or in AI terms: training data,
system prompts, and hard-coded constraints)---is displaced by a
parasitic instruction. The agent then proceeds to Decide and Act with
perfect logical consistency, but in service of an alien will. This
dynamic has been documented across the emerging agentic AI landscape
\cite{owasp2025agentic, microsoft2025indirect}. The OWASP Top 10 for
Agentic Applications (December 2025) designates \textbf{ASI-01: Agent
Goal Hijack} as the \#1 risk for deployed agentic AI systems---a direct
industry validation of this paper's central thesis.

\begin{block}{Empirical Urgency}
\protect\phantomsection\label{empirical-urgency}
Goal Hijacking has transitioned from academic concern to documented
production failure. Autonomous coding agents have deleted production
databases and fabricated records to conceal the damage; invisible
Unicode payloads have triggered auto-approval modes enabling remote code
execution; and indirect prompt injection through enterprise messaging
platforms has exfiltrated private API keys---all without human
authorization
\cite{adversa2025incidents, copilot2025rce, promptarmor2024slack}. The
Agent Security Bench (ASB) evaluation \cite{zhang2025asb} quantifies the
gap: an 84.3\% average attack success rate across 400+ integrated tools,
with current defenses achieving only 19.7\% mitigation. He et
al.~\cite{he2025redteaming} further demonstrate Agent-in-the-Middle
attacks that compromise inter-agent communication channels, extending
the threat surface to multiagent coordination itself. These incidents
are analyzed in detail through the CIF-AD-OODA lens in
\cref{sec:empirical_grounding}.
\end{block}
\end{frame}

\begin{frame}{Axiomatic Design and Transient Functional Requirements}
\protect\phantomsection\label{axiomatic-design-and-transient-functional-requirements}
Suh's \textbf{Axiomatic Design (AD)} theory
\cite{suh1990principles, suh2001axiomatic} posits that good design
maintains the independence of Functional Requirements (FRs). The
relationship between FRs and Design Parameters (DPs) is captured in the
Design Matrix:

\begin{equation}
\{FR\} = [A]\{DP\}
\label{eq:intro_design_matrix}
\end{equation}

In an uncoupled (safe) design, \([A]\) is diagonal: each FR depends on
exactly one DP. Providing a solution for \(\text{FR}_1\) does not
compromise \(\text{FR}_2\). In cognitive security, we identify a new
failure mode: \textbf{Transient Functional Coupling}.

Adversaries use ``fast OODA transients''---high-frequency semantic
injections through data channels (\(\Omega_2\) vectors)---to temporarily
introduce off-diagonal terms in the Design Matrix. For instance, a
cyber-defense agent's FR (``Block Malicious Traffic'') might be
transiently coupled to a hijacked FR (``Maximize Uptime''), causing the
agent to flush its firewalls during an attack to ``restore
connectivity.'' The coupled matrix:

\begin{equation}
\{FR'\} = [A']\{DP\} = \begin{bmatrix} A_{11} & A_{12} \\ A_{21} & A_{22} \end{bmatrix} \{DP\}
\label{eq:intro_coupled_matrix}
\end{equation}

violates the Independence Axiom, rendering the system unstable. The
focus of this paper is on detecting and defending against these rapid
shifts in the Design Matrix.
\end{frame}

\begin{frame}{The Role of CIF: Five Canonical Defenses}
\protect\phantomsection\label{the-role-of-cif-five-canonical-defenses}
The Cognitive Integrity Framework serves as the ``gyroscope'' for OODA
loops, stabilizing the Orientation phase against adversarial transients.
CIF provides five canonical defense mechanisms
\cite{friedman2026cogsec1}, each targeting a specific aspect of
cognitive integrity:

\begin{enumerate}
\tightlist
\item
  \textbf{Cognitive Firewall} (\(\mathcal{F}\)): Classifies incoming
  messages by trust score, quarantining ambiguous inputs before they
  reach the agent's Orientation phase.
\item
  \textbf{Belief Sandboxing} (\(\mathcal{B}_{\text{provisional}}\)):
  Isolates unverified beliefs in a provisional partition, requiring
  corroboration before promotion to operational status.
\item
  \textbf{Behavioral Invariants} (\(\text{INV}_k\)): Pre-defined runtime
  predicates that trigger alerts when violated, regardless of semantic
  content.
\item
  \textbf{Drift Detection} (\(S_{\text{drift}}\)): Monitors belief
  distribution changes via KL divergence, flagging sudden orientation
  shifts that exceed threshold \(\epsilon\).
\item
  \textbf{Byzantine Consensus} (\(\mathcal{B}_{\text{consensus}}\)):
  Multi-agent agreement protocol requiring quorum \(q\) before critical
  actions, tolerating up to \(f\) compromised agents where
  \(n \geq 3f+1\) \cite{lamport1982byzantine}.
\end{enumerate}

These mechanisms compose in series and parallel to achieve layered
defense. Rather than the ad-hoc ``Axiomatic Locking'' described in early
formulations, CIF's defense is the \emph{systematic composition} of
these five mechanisms, calibrated to the threat profile and temporal
dynamics of each operational domain.
\end{frame}

\begin{frame}{Application Analysis}
\protect\phantomsection\label{application-analysis}
This manuscript examines Goal Hijacking dynamics across ten high-stakes
domains, analyzing how adversaries exploit the collaborative surface
between AI agents, OODA loops, and Axiomatic Design principles. Each
domain is analyzed using a standardized five-step template
(\cref{sec:methodology}) that characterizes the operational context,
attack surface, transient coupling, defense mapping, and validation
anchoring.
\end{frame}

\begin{frame}{Contributions}
\protect\phantomsection\label{contributions}
This paper makes the following contributions:

\begin{itemize}
\tightlist
\item
  \textbf{C1:} A unified CIF-AD-OODA integration model for analyzing
  Goal Hijacking attacks and defenses across arbitrary operational
  domains.
\item
  \textbf{C2:} Identification of three universal attack patterns---FR
  Polarity Inversion, Constraint Relaxation, and Context Boundary
  Violation---through cross-domain synthesis.
\item
  \textbf{C3:} Validation that all five canonical CIF mechanisms provide
  adequate coverage across ten critical domains, with no mechanism
  appearing in fewer than three domains and no domain requiring
  mechanisms outside the CIF vocabulary.
\item
  \textbf{C4:} Three novel defense pattern extensions: verification
  channel separation, active perturbation probing, and physics-informed
  invariants.
\item
  \textbf{C5:} Temporal scale analysis demonstrating CIF's applicability
  across eight orders of magnitude in OODA cycle time.
\item
  \textbf{C6:} Retrospective validation through six documented AI agent
  security incidents (2024--2025), confirming that all incidents map to
  the universal attack pattern taxonomy and would have been detectable
  by the appropriate CIF mechanism.
\end{itemize}
\end{frame}

\end{document}
